% =====================================================================
% PAPER III — Eisenstein–Picard by verified finite elements
% Level-1 gap for PSL(2,Z[omega]) and first congruence rung Gamma_0(1-omega)
% ---------------------------------------------------------------------
% Sources in-repo:
%   eisenstein_numbers/PROOF.md, PAIRING_MATRICES.md, T0_AREA.md,
%   GEOMETRY.md, CERT_OMEGA_P.md, cert_omega.py, cert_omega_p.py,
%   pairing_matrices.py, non_neumann_omega.py
% Parallel method paper: papers/paper1_hammer.tex (Picard level 1)
% =====================================================================
\documentclass[11pt]{amsart}
\usepackage{amsmath,amssymb,amsthm}
\newtheorem{theorem}{Theorem}[section]
\newtheorem{lemma}[theorem]{Lemma}
\newtheorem{proposition}[theorem]{Proposition}
\newtheorem{remark}[theorem]{Remark}
\newcommand{\HH}{\mathbb{H}^3}
\newcommand{\ZZ}{\mathbb{Z}}
\newcommand{\CC}{\mathbb{C}}
\newcommand{\PSL}{\mathrm{PSL}}
\newcommand{\note}[1]{\par\smallskip\noindent\textbf{[[#1]]}\par\smallskip}

\title[Verified exclusion for the Eisenstein--Picard orbifold]{Exceptional
eigenvalues of Bianchi orbifolds by verified finite elements.\\
III: The Eisenstein--Picard orbifold and its first congruence cover}
\author{}
\date{Draft, 2026-07-12}

\begin{document}
\begin{abstract}
We apply the verified finite-element exclusion method of Paper~I to the
Eisenstein--Picard group $\Gamma=\PSL(2,\ZZ[\omega])$, $\omega=e^{2\pi i/3}$.
Using the Elstrodt--Grunewald--Mennicke fundamental domain, Crouzeix--Raviart
elements on a truncated core, ball-arithmetic enclosures, and Rump's
positive-definiteness certificate, we prove that
$L^2(\Gamma\backslash\HH)$ has no Laplace eigenvalue in $(0,1)$.  The same
architecture, with multi-copy reference cells and two-cusp weights, yields
the first congruence rung: $\Gamma_0(1-\omega)\backslash\HH$ likewise has no
eigenvalue in $(0,1)$.  The method is independent of the Selberg trace
formula.
\end{abstract}
\maketitle

\section{Introduction}
\note{PROOF.md Theorems A--B; parent README context.}
\begin{itemize}
\item Bianchi groups for $d=-3$; comparison with Picard $d=-1$ (Paper~I)
  and the Gaussian congruence ladder (Paper~II / in-repo Thms~1--4).
\item Main theorems (level~1 and $\Gamma_0(1-\omega)$).
\item Independence from the trace formula; relation to the in-repo
  STF/CE gate $B(\mathbb{Q}(\omega))<1$.
\item Reproducibility: \texttt{eisenstein\_numbers/cert\_omega.py},
  \texttt{cert\_omega\_p.py}.
\end{itemize}

\begin{theorem}[Level 1]\label{thm:A}
The Laplace--Beltrami operator on
$L^2(\PSL(2,\ZZ[\omega])\backslash\HH)$ has no eigenvalue in $(0,1)$.
Combined with the spectral decomposition of cofinite Kleinian groups,
$\lambda_1\ge 1$.
\end{theorem}

\begin{theorem}[First congruence rung]\label{thm:B}
Let $\mathfrak{p}=(1-\omega)\subset\ZZ[\omega]$ and
$\Gamma_0(\mathfrak{p})\subset\PSL(2,\ZZ[\omega])$.  Then
$L^2(\Gamma_0(\mathfrak{p})\backslash\HH)$ has no Laplace eigenvalue in
$(0,1)$.
\end{theorem}

\section{Geometry}
\note{GEOMETRY.md, geometry\_fund.py, EGM98 §7.3, DP20 §2.3.}
Fundamental domain $F_3=B_3\cap(P_3\times\mathbb{R}_{>0})$,
$|T|=\mathrm{area}(P_3)=\sqrt3/6$, truncated core
$K_Y=F_3\cap\{t\le Y\}$ with $Y=1.25$.
\begin{equation*}
  \mathrm{vol}(K_Y)=\mathrm{vol}(F)-|T|/(2Y^2).
\end{equation*}

\subsection*{Pairing matrices}
\note{PAIRING\_MATRICES.md, GENERATOR\_CYCLES.md.}
Generators in $\mathrm{SL}(2,\ZZ[\omega])$:
\begin{align*}
  T_1&=\begin{pmatrix}1&1\\0&1\end{pmatrix},&
  T_\omega&=\begin{pmatrix}1&\omega\\0&1\end{pmatrix},\\
  U&=\begin{pmatrix}\omega^2&0\\0&\omega\end{pmatrix},&
  S&=\begin{pmatrix}0&-1\\1&0\end{pmatrix}.
\end{align*}
Reference face maps on $P_3$: $T_1$ pairs RIGHT$\leftrightarrow$LEFT,
$T_\omega$ pairs LOW$\leftrightarrow$UP, $U$ acts by $z\mapsto\omega z$ on
vertical walls, $S$ self-pairs the floor.  Ideal-vertex cycles of inverses
are recorded and machine-checked against a fixed $\Gamma_\infty$-section.

\begin{remark}[Neumann vs non-Neumann]
The certified level-1 FE space uses Neumann free $H^1(K)$ (relaxation:
larger space).  The optional module \texttt{non\_neumann\_omega.py}
imposes the pairing dictionary on CR face dofs; positivity there is
implied by the Neumann certificate and is checked in floating point for
both $\kappa_{\mathrm{sc}}$ (Lemma~I1) and the sharper CZZ constant.
\end{remark}

\section{Analytic criterion}
\note{DESIGN.md (Picard) with $|T|=\sqrt3/6$; T0\_AREA.md for two cusps.}
Lax--Phillips reduction on the truncated core; single-cusp weight
$\beta(s)=(1-s)/(|T|Y^2)$ at level~1.  For $\Gamma_0(\mathfrak{p})$ with
$N(\mathfrak{p})=q$ prime (field residue), two cusps $\infty,0$ with
\begin{equation*}
  |T_\infty|=\sqrt3/6,\qquad |T_0|=q\cdot|T_\infty|,
\end{equation*}
and
\begin{equation*}
  \beta_\infty(s)=\frac{1-s}{|T_\infty|Y^2},\qquad
  \beta_0(s)=\frac{1-s}{|T_0|Y^2}.
\end{equation*}

\section{Guaranteed FEM lower bound}
\note{lower\_bound\_theory.md Theorem G1; CERT\_OMEGA\_P.md Theorem G1p.}
Crouzeix--Raviart on a polyhedral inner mesh of $K_Y$ (Lemma~G floor lift).
Finite checks per $s$-window: $c_e>d_e$ and $N_h-D_h\succeq 0$.
Congruence: Reference-Cell Principle with
$\mathrm{NC}=N(\mathfrak{p})+1$ isometric copies; Lemma~D0
($\sigma_h=\sqrt{\mathrm{index}}\,\alpha_h^{\mathrm{ref}}$) when multi-copy
tops remain free.

\section{Verified certificate}
\note{cert\_omega.py, cert\_omega\_p.py, Rump BIT 46.}
\begin{itemize}
\item Geometry and element matrices in Arb (python-flint, 128-bit).
\item Scalar window inequalities by rigorous ball comparison.
\item Matrix positive-semidefiniteness by Rump (BIT~46, 2006) with
  power-of-two diagonal equilibration and per-row radii.
\end{itemize}

\subsection*{Frozen parameters (Theorem~\ref{thm:A})}
Mesh $P_3$, $N_{\mathrm{tri}}=6$, $N_3=3$ (648 tets, 1440 CR dofs);
$Y=1.25$, $\rho=55$, $\theta=0.7$, $\theta_2=\theta'=\alpha=0.5$,
$\nu^*=1.05$; eight uniform windows on $[0,1]$;
$\kappa_{\mathrm{sc}}=\sqrt{1/\pi^2+1/15}$.

\subsection*{Frozen parameters (Theorem~\ref{thm:B})}
Same reference mesh; $\mathrm{NC}=4$; $\theta=0.5$, $\theta_2=0.9$,
$\alpha=0.2$, $\theta_4=0.5$, $\tilde\rho=9$; $n\approx 5404$ dofs.

\section{Reproduce}
\begin{verbatim}
cd eisenstein_numbers
python -u geometry_fund.py
python -u pairing_matrices.py
python -u check_generators.py
python -u non_neumann_omega.py      # optional paired FE space
python -u cert_omega.py 6 3         # Theorem A
python -u cert_omega_p.py 3 6 3     # Theorem B
\end{verbatim}

\section{External citations (not re-proved)}
\begin{enumerate}
\item Friedman, arXiv:math/0612807, Thm~3.8.1 (spectral decomposition).
\item Elstrodt--Grunewald--Mennicke, \emph{Groups Acting on Hyperbolic Space},
  §7.2--7.3; planar domain $P_3$.
\item Shimizu / Leutbecher (precisely invariant horoball).
\item Morrey--Nirenberg (analyticity / unique continuation).
\item Rump, BIT Numer.\ Math.\ 46 (2006).
\item Payne--Weinberger / Bebendorf (CR constant; optional CZZ20/CP22).
\end{enumerate}

\section{What is not claimed}
No statement for $N(\mathfrak{p})\in\{7,13\}$ beyond combinatorics;
no re-proof of the Gaussian ladder; no use of the Selberg trace formula
in the certificates above.

\begin{thebibliography}{9}
\bibitem{EGM98} J.~Elstrodt, F.~Grunewald, J.~Mennicke,
  \emph{Groups Acting on Hyperbolic Space}, Springer, 1998.
\bibitem{Rump06} S.M.~Rump, Verification of positive definiteness,
  BIT Numer.\ Math.\ 46 (2006), 433--452.
\bibitem{Friedman} J.~Friedman, The Selberg trace formula for
  $\mathrm{PSL}(2,O_K)$, arXiv:math/0612807.
\end{thebibliography}

\end{document}
